\chapter{Внутренняя геометрия и алгебра тензоров}

Тензорная нотация полезна лишь тогда, когда часть размышлений она выполняет за нас. Хорошо устроенная нотация ведет себя как набор механических деталей: верхний индекс сочетается с нижним, несовместимые объекты не позволяют себя складывать, а допустимая свертка оставляет объект, закон преобразования которого можно прочитать по оставшимся индексам. Непосредственная цель этого аппарата — геометрическая. Мы хотим отличать подлинную кривизну от кривизны, возникающей лишь потому, что координаты были искривлены.

\section{Кривизна изнутри}

Начнем с обычного листа бумаги. Его можно свернуть в цилиндр без растяжения и разрыва. Для наблюдателя в окружающей комнате лист явно изогнут; тем не менее расстояния, измеряемые на нем, не меняются. Его внешняя кривизна изменилась, а внутренняя геометрия — нет.

Представим теперь маленькое существо, заключенное на этом листе. Оно не может видеть окружающую комнату. Оно может измерять длины, сравнивать соседние направления, строить треугольники и переносить мерные стержни с места на место. Именно таких внутренних экспериментов геометрия и должна требовать. Соответствующий вопрос в теории тяготения состоит в том, представляет ли поле истинную кривизну пространства-времени или лишь криволинейное описание плоского пространства-времени.

Один из способов сделать эту внутреннюю точку зрения наглядной — заменить поверхность сетью точек и звеньев. Соединим соседние точки, запишем длину каждого звена и допустим, что соединения шарнирны. Равносторонняя треугольная решетка лежит в плоскости. Сгибание шарниров не изменяет закодированную геометрию, потому что ее всегда можно вернуть в плоскость без изменения длин.

Теперь удлиним шесть звеньев, сходящихся в одной вершине, оставив окружающие звенья фиксированными. Центральная вершина уже не может оставаться в плоскости, не нарушая предписанных длин. Она должна подняться, образовав выпуклость. Чтобы распознать эту кривизну, не понадобился никакой взгляд извне: она уже содержалась в несовместимых расстояниях между соседними точками.

\begin{figure}[H]
\centering
\includegraphics[width=\linewidth]{lecture_02_figure_01.png}
\caption{Треугольная модель внутренней геометрии из звеньев. Изменение звеньев вокруг одной вершины приводит к выпуклости, которую нельзя сделать плоской без изменения длин.}
\lectureframe{00:06:40}
\end{figure}

\begin{classroomqa}
\audiencequestion{Каково же тогда точное определение плоского пространства?}
\lecturerresponse{Операциональный ответ таков: его внутренние соотношения расстояний можно реализовать без искажений в евклидовом расположении. В языке, который будет развит ниже, риманово пространство плоско, если во всей интересующей нас области можно найти координаты, в которых метрика является постоянной дельтой Кронекера. В каждой отдельной точке этого всегда можно добиться локально; добиться этого на целой области — нетривиальное условие. Тензор кривизны в конечном счете сформулирует это условие без обращения к специальному выбору координат.
\lecturetimestamp{00:09:15}}
\end{classroomqa}

В этом ответе мы использовали метрику прежде, чем успели ее ввести. Поэтому мы временно перейдем от кривизны к тензорному анализу. Этот обход не факультативен: метрику нужно распознать как тензор, прежде чем можно будет отделить ее координатную зависимость от подлинной геометрии.

\section{Поля и координатные описания}

Тензор в общей теории относительности обычно является тензорным \emph{полем}: в каждой точке имеется тензор, и его компоненты меняются от точки к точке. Скалярное поле — простейший пример. Температуру на поверхности можно представить числом \(S(P)\) в каждой точке \(P\). Два человека могут обозначить одну и ту же точку разными координатами, но относительно температуры в ней они должны согласиться.

\begin{figure}[H]
\centering
\includegraphics[width=\linewidth]{lecture_02_figure_02.png}
\caption{Скалярное поле и два координатных обозначения одних и тех же точек.}
\lectureframe{00:11:25}
\end{figure}

Пусть \(x^m\) и \(y^m\) — две гладкие системы координат. Каждая из них есть отображение точек в упорядоченные наборы чисел, и в общей координатной области эти два набора связаны взаимно обратными функциями,
\begin{equation}
 y^m=y^m(x^1,\ldots,x^d),
 \qquad
 x^m=x^m(y^1,\ldots,y^d).
\label{eq:coordinate-maps}
\end{equation}
Гладкость позволяет дифференцировать, а обратимость гарантирует, что ни одна точка не теряется и не отождествляется с другой лишь из-за смены координат.

\begin{figure}[H]
\centering
\includegraphics[width=\linewidth]{lecture_02_figure_03.png}
\caption{Два гладких обратимых координатных отображения, описывающих одну и ту же геометрическую область.}
\lectureframe{00:13:10}
\end{figure}

Для скалярного поля закон преобразования просто таков
\begin{equation}
 S'(y)=S(x).
\label{eq:scalar-law}
\end{equation}
Функции по обе стороны могут выглядеть по-разному, поскольку их аргументы различны, но в соответствующих точках их значения совпадают. Это прототип инвариантного утверждения.

\section{Косоугольные координаты и два вида компонент}

Прежде чем определять общие тензоры через законы преобразования, полезно понять, почему верхние и нижние индексы вектора нельзя считать просто декоративными вариантами друг друга.

Рассмотрим координатную сетку, оси которой не являются ни взаимно перпендикулярными, ни нормированными. В данной точке увеличение \(x^m\) на одну координатную единицу задает базисный вектор \(\mathbf e_m\). На языке координат,
\begin{equation}
 \mathbf e_m=\frac{\partial \mathbf r}{\partial x^m}.
\label{eq:coordinate-basis}
\end{equation}
Векторы \(\mathbf e_m\) не обязаны иметь единичную длину и не обязаны быть ортогональными.

\begin{figure}[H]
\centering
\includegraphics[width=\linewidth]{lecture_02_figure_04.png}
\caption{Косоугольная координатная сетка и ее координатные базисные векторы.}
\lectureframe{00:17:40}
\end{figure}

Любой вектор \(\mathbf V\) можно разложить по этому базису:
\begin{equation}
 \mathbf V=V^m\mathbf e_m.
\label{eq:contravariant-expansion}
\end{equation}
Числа \(V^m\) — это \emph{контравариантные компоненты}. Это коэффициенты, необходимые для того, чтобы собрать геометрический вектор из базисных векторов.

\begin{figure}[H]
\centering
\includegraphics[width=\linewidth]{lecture_02_figure_05.png}
\caption{Разложение вектора по неортогональному координатному базису.}
\lectureframe{00:19:20}
\end{figure}

Существует и другой естественный набор чисел:
\begin{equation}
 V_m=\mathbf V\mathbin{\boldsymbol{\cdot}}\mathbf e_m.
\label{eq:covariant-projection}
\end{equation}
Это \emph{ковариантные компоненты}. В косоугольном базисе это не просто те же коэффициенты с перенесенным индексом. Подставляя \eqref{eq:contravariant-expansion}, получаем
\begin{equation}
 V_m
 =V^n\bigl(\mathbf e_n\mathbin{\boldsymbol{\cdot}}\mathbf e_m\bigr).
\label{eq:covariant-from-basis}
\end{equation}

\begin{figure}[H]
\centering
\includegraphics[width=\linewidth]{lecture_02_figure_06.png}
\caption{Ковариантные компоненты, полученные взятием скалярных произведений с координатным базисом.}
\lectureframe{00:21:10}
\end{figure}

Скалярные произведения базисных векторов образуют матрицу Грама
\begin{equation}
 g_{mn}=\mathbf e_m\mathbin{\boldsymbol{\cdot}}\mathbf e_n.
\label{eq:metric-gram}
\end{equation}
Следовательно,
\begin{equation}
 V_m=g_{mn}V^n.
\label{eq:lower-vector}
\end{equation}
Это первое появление метрики. Она превращает коэффициенты разложения в проекции, поскольку содержит как длины, так и взаимные углы координатного базиса.

\begin{figure}[H]
\centering
\includegraphics[width=\linewidth]{lecture_02_figure_07.png}
\caption{Метрика как матрица Грама базисных векторов.}
\lectureframe{00:23:40}
\end{figure}

Та же матрица вычисляет длину:
\begin{align}
 \mathbf V\mathbin{\boldsymbol{\cdot}}\mathbf V
 &=(V^m\mathbf e_m)\mathbin{\boldsymbol{\cdot}}
   (V^n\mathbf e_n) \notag\\
 &=g_{mn}V^mV^n.
\label{eq:vector-length}
\end{align}

\begin{figure}[H]
\centering
\includegraphics[width=\linewidth]{lecture_02_figure_08.png}
\caption{Норма вектора, записанная через метрику и контравариантные компоненты.}
\lectureframe{00:25:00}
\end{figure}

\begin{classroomqa}
\audiencequestion{Зачем вводить эту геометрическую конструкцию до формальных законов преобразования?}
\lecturerresponse{Она придает индексам конкретный смысл. Верхние компоненты — это коэффициенты разложения, нижние компоненты — линейные измерения вектора, а метрика связывает одно с другим. Тогда формальное определение скажет нам, что эти роли сохраняются при произвольных гладких заменах координат. \lecturetimestamp{00:26:07}}
\end{classroomqa}

\begin{classroomqa}
\audiencequestion{Когда ковариантные и контравариантные компоненты имеют одинаковые численные значения?}
\lecturerresponse{В ортонормированном евклидовом базисе \(g_{mn}=\delta_{mn}\), так что \(V_m=V^m\) покомпонентно. Одной лишь ортогональности недостаточно: если базисные векторы имеют неединичные длины, диагональная метрика все равно вносит масштабные множители.
\lecturetimestamp{00:26:19}}
\end{classroomqa}

Полярные координаты дают наглядное предостережение. Радиальные и угловые координатные линии пересекаются ортогонально, однако фиксированному изменению \(d\theta\) соответствует физическое расстояние \(r\,d\theta\). На плоскости,
\begin{equation}
 ds^2=dr^2+r^2d\theta^2,
\label{eq:polar-metric}
\end{equation}
поэтому базис ортогонален, но не декартов.

\begin{figure}[H]
\centering
\includegraphics[width=\linewidth]{lecture_02_figure_09.png}
\caption{Полярные координаты: ортогональные координатные направления с угловым масштабным множителем, растущим с радиусом.}
\lectureframe{00:28:10}
\end{figure}

\begin{classroomqa}
\audiencequestion{Являются ли ковариантные векторы действительно иными векторами или это лишь второй способ задания того же вектора?}
\lecturerresponse{До задания метрики вектор и ковектор — разные виды линейных объектов. Вектор принадлежит касательному пространству; ковектор есть линейный функционал на этом пространстве. Невырожденная метрика устанавливает между ними взаимно однозначное соответствие, задаваемое \(V_m=g_{mn}V^n\), но это соответствие не следует принимать за тождество. \lecturetimestamp{00:29:10}}
\end{classroomqa}

\section{Законы преобразования определяют тензор}

Самое надежное формальное определение — операциональное: массив компонент является тензором, если его компоненты преобразуются с соответствующими множителями Якоби. Дифференцируя \eqref{eq:coordinate-maps}, получаем
\begin{equation}
 dy^m=\frac{\partial y^m}{\partial x^n}\,dx^n.
\label{eq:differential-law}
\end{equation}
Любой контравариантный вектор подчиняется тому же правилу,
\begin{equation}
 V'^m(y)
 =\frac{\partial y^m}{\partial x^n}V^n(x).
\label{eq:contravariant-law}
\end{equation}

\begin{figure}[H]
\centering
\includegraphics[width=\linewidth]{lecture_02_figure_10.png}
\caption{Контравариантный закон преобразования, построенный по образцу преобразования координатных дифференциалов.}
\lectureframe{00:32:50}
\end{figure}

Ковектор преобразуется с помощью обратного якобиана:
\begin{equation}
 W'_m(y)
 =\frac{\partial x^n}{\partial y^m}W_n(x).
\label{eq:covariant-law}
\end{equation}
Градиент скаляра — стандартный пример. Применяя правило цепочки к \eqref{eq:scalar-law},
\begin{equation}
 \frac{\partial S}{\partial y^m}
 =\frac{\partial x^n}{\partial y^m}
  \frac{\partial S}{\partial x^n}.
\label{eq:gradient-law}
\end{equation}

\begin{figure}[H]
\centering
\includegraphics[width=\linewidth]{lecture_02_figure_11.png}
\caption{Градиент скаляра, который по правилу цепочки преобразуется как ковектор.}
\lectureframe{00:34:20}
\end{figure}

\begin{classroomqa}
\audiencequestion{Обязательно ли верхние и нижние компоненты имеют одни и те же физические единицы?}
\lecturerresponse{Нет. Координатные компоненты наследуют единицы измерения и соглашения о масштабе от координат и их базисных векторов. Если одна координата — длина, а другая — угол, их компоненты могут даже не иметь одинаковой размерности. Полная свёртка — это объект, физическая размерность которого не зависит от координат.
\lecturetimestamp{00:35:41}}
\end{classroomqa}

\begin{classroomqa}
\audiencequestion{Почему здесь производные частные, а не обычные?}
\lecturerresponse{Поле зависит от нескольких координат. Якобиан показывает, как изменяется одна координатная функция, когда меняется один аргумент, а остальные удерживаются фиксированными, поэтому его элементы — частные производные.
\lecturetimestamp{00:36:45}}
\end{classroomqa}

Теперь эта схема механически продолжается. Каждый верхний штрихованный индекс вносит один множитель \(\partial y/\partial x\); каждый нижний штрихованный индекс вносит один множитель \(\partial x/\partial y\). Смешанный тензор преобразуется по закону
\begin{equation}
 T'^m{}_n
 =
 \frac{\partial y^m}{\partial x^p}
 \frac{\partial x^q}{\partial y^n}
 T^p{}_q,
\label{eq:mixed-law}
\end{equation}
а ковариантный тензор второго ранга — по закону
\begin{equation}
 W'_{mn}
 =
 \frac{\partial x^p}{\partial y^m}
 \frac{\partial x^q}{\partial y^n}
 W_{pq}.
\label{eq:rank-two-covariant-law}
\end{equation}

\begin{figure}[H]
\centering
\includegraphics[width=\linewidth]{lecture_02_figure_12.png}
\caption{Закон преобразования тензора с одним верхним и одним нижним индексом.}
\lectureframe{00:40:45}
\end{figure}

\begin{figure}[H]
\centering
\includegraphics[width=\linewidth]{lecture_02_figure_13.png}
\caption{Ковариантный закон преобразования второго ранга: по одному обратному якобиану на каждый нижний индекс.}
\lectureframe{00:42:15}
\end{figure}

Повторяющиеся индексы суммируются, если не оговорено иное. Это соглашение устраняет знаки суммирования, сохраняя всю информацию, необходимую, чтобы их восстановить. Повторяющийся немой индекс должен встречаться в слагаемом ровно дважды: один раз сверху и один раз снизу; свободный индекс встречается по одному разу в каждом члене уравнения.

\begin{classroomqa}
\audiencequestion{Кто ввёл это обозначение суммирования?}
\lecturerresponse{Обычно его приписывают Эйнштейну, поэтому оно и называется соглашением Эйнштейна о суммировании. Его ценность не сводится к историческому украшению: рисунок индексов позволяет с первого взгляда увидеть допустимые тензорные операции. \lecturetimestamp{00:43:21}}
\end{classroomqa}

\section{Тензорные уравнения и истина, не зависящая от выбора координат}

Компоненты меняются при смене координат, но тензорное уравнение остаётся истинным в любой системе координат. Если два тензора одного и того же типа удовлетворяют \(W=T\), то их разность есть нулевой тензор:
\begin{equation}
 W^{m\cdots}{}_{\cdots n}
 -T^{m\cdots}{}_{\cdots n}=0.
\label{eq:tensor-equality}
\end{equation}
Закон преобразования линеен, поэтому нулевые компоненты в одной системе дают нулевые компоненты во всякой системе. Совпадающие свободные индексы в левой и правой частях необходимы. Уравнение между объектами разных тензорных типов не имеет координатно-независимого смысла.

\begin{figure}[H]
\centering
\includegraphics[width=\linewidth]{lecture_02_figure_14.png}
\caption{Совпадающая индексная структура превращает равенство компонент в тензорное уравнение, верное в любой системе координат.}
\lectureframe{00:46:35}
\end{figure}

\begin{classroomqa}
\audiencequestion{Что именно инвариантно, если компоненты вектора меняются?}
\lecturerresponse{Инвариантны геометрический тензор и истинность тензорного уравнения; список его компонент в координатах — нет. Скаляр — частный случай без индексов, поэтому его численное значение в точке одинаково в любой системе координат. \lecturetimestamp{00:47:19}}
\end{classroomqa}

\begin{classroomqa}
\audiencequestion{Является ли норма вектора соответствующим инвариантом, и означает ли нулевая норма, что вектор равен нулю?}
\lecturerresponse{Когда задана метрика, свёртка \(g_{mn}V^mV^n\) является скаляром и, следовательно, инвариантным квадратом нормы. В положительно определённой геометрии она обращается в нуль только для нулевого вектора. В лоренцевом пространстве-времени ненулевой светоподобный вектор может иметь нулевую норму, так что эти два утверждения нужно различать. \lecturetimestamp{00:48:04}}
\end{classroomqa}

\section{Допустимые алгебраические операции над тензорами}

На доске собраны основные операции, которые нам понадобятся: умножение на скаляр, сложение и вычитание тензоров одного типа, образование тензорных произведений, свёртка верхнего индекса с нижним и, позднее, ковариантное дифференцирование.

\begin{figure}[H]
\centering
\includegraphics[width=\linewidth]{lecture_02_figure_15.png}
\caption{Перечень тензорных операций на доске: сложение, умножение, свёртка и дифференцирование.}
\lectureframe{00:50:40}
\end{figure}

\begin{classroomqa}
\audiencequestion{Дифференцирование по чему?}
\lecturerresponse{По положению, а значит, по координатам. Обычные производные компонент, вообще говоря, не преобразуются как тензоры, потому что сами якобианы меняются. Ковариантное дифференцирование исправляет это, но его построение отложено до тех пор, пока не будут подготовлены геометрические предварительные понятия. \lecturetimestamp{00:51:49}}
\end{classroomqa}

Умножение на скаляр не вызывает затруднений. Сложение и вычитание требуют совпадения тензорного типа:
\begin{equation}
 C^m{}_n=A^m{}_n+B^m{}_n,
\qquad
 D^m{}_n=A^m{}_n-B^m{}_n.
\label{eq:tensor-addition}
\end{equation}
Нельзя складывать вектор с ковектором или тензор второго ранга со скаляром, так же как нельзя складывать скорость с площадью.

\begin{figure}[H]
\centering
\includegraphics[width=\linewidth]{lecture_02_figure_16.png}
\caption{Сложение покомпонентно определено только для тензоров с одинаковой схемой свободных индексов.}
\lectureframe{00:53:00}
\end{figure}

Тензорное произведение объединяет все индексы своих множителей. Например,
\begin{equation}
 T^{mn}=V^mW^n,
\qquad
 S^m{}_n=V^mU_n.
\label{eq:tensor-products}
\end{equation}
В четырёх измерениях \(T^{mn}\) имеет 16 компонент. Ни по одному индексу суммирование не выполнено, так что это внешнее произведение, а не скалярное произведение.

\begin{figure}[H]
\centering
\includegraphics[width=\linewidth]{lecture_02_figure_17.png}
\caption{Тензорное произведение, сохраняющее по одному свободному индексу от каждого векторного множителя.}
\lectureframe{00:55:15}
\end{figure}

\begin{classroomqa}
\audiencequestion{Является ли умножение двух векторов здесь скалярным произведением?}
\lecturerresponse{Нет. Скалярное произведение свёртывает индекс и даёт скаляр. Тензорное произведение оставляет оба индекса свободными и потому даёт тензор второго ранга. В четырёх измерениях это 16 компонент, а не одна. \lecturetimestamp{00:56:36}}
\end{classroomqa}

\begin{classroomqa}
\audiencequestion{Поскольку числовые компоненты коммутируют, не означает ли это, что перестановка двух векторных множителей ничего не меняет?}
\lecturerresponse{Для обычных скалярных компонент \(V^mW^n=W^nV^m\). Но тензорные места упорядочены. Если переставить множители, а затем назвать первое место \(m\), а второе \(n\), получится транспонированный массив; \(T^{mn}\) не обязано равняться \(T^{nm}\).
\lecturetimestamp{00:58:52}}
\end{classroomqa}

\section{Свёртка и лемма об обратном якобиане}

Ключ к свёртке — правило цепочки. Поскольку координатные отображения обратны друг другу,
\begin{equation}
 \frac{\partial x^b}{\partial y^m}
 \frac{\partial y^m}{\partial x^a}
 =
 \frac{\partial x^b}{\partial x^a}
 =\delta^b{}_a.
\label{eq:inverse-jacobian}
\end{equation}
эти две матрицы Якоби взаимно обратны, и их произведение есть дельта Кронекера.

\begin{figure}[H]
\centering
\includegraphics[width=\linewidth]{lecture_02_figure_18.png}
\caption{Тождество для обратного якобиана, подготовленное как лемма, лежащая в основе тензорной свёртки.}
\lectureframe{01:02:15}
\end{figure}

\begin{classroomqa}
\audiencequestion{Что это за объект — обратный якобиан?}
\lecturerresponse{Это матрица производных обратного координатного отображения. Её элементы — \(\partial x^b/\partial y^m\), и при матричном умножении на \(\partial y^m/\partial x^a\) по многомерному правилу цепочки получается единичная матрица. \lecturetimestamp{01:02:37}}
\end{classroomqa}

Короткое доказательство проясняет обозначения. Для любой скалярной функции \(f\)
\begin{equation}
 \frac{\partial f}{\partial y^m}
 \frac{\partial y^m}{\partial x^a}
 =\frac{\partial f}{\partial x^a}.
\end{equation}
Выбрав \(f=x^b\), получаем
\begin{equation}
 \frac{\partial x^b}{\partial y^m}
 \frac{\partial y^m}{\partial x^a}
 =\frac{\partial x^b}{\partial x^a}
 =\delta^b{}_a.
\end{equation}

\begin{figure}[H]
\centering
\includegraphics[width=\linewidth]{lecture_02_figure_19.png}
\caption{Правило цепочки, сводящее обратные якобианы к дельте Кронекера.}
\lectureframe{01:04:55}
\end{figure}

\begin{classroomqa}
\audiencequestion{Представляет ли \(\delta^b{}_a\) тождественное отображение только в декартовых координатах?}
\lecturerresponse{Нет. Смешанный тензор \(\delta^b{}_a\) представляет тождественное отображение в любой системе координат. Евклидова метрика, записанная как \(\delta_{ab}\), характерна только для ортонормированных декартовых координат, но этот тождественный тензор со смешанными индексами не следует путать с той метрикой.
\lecturetimestamp{01:04:58}}
\end{classroomqa}

Теперь рассмотрим смешанный тензор \(T^m{}_n\). После преобразования его след равен
\begin{align}
 T'^m{}_m
 &=
 \frac{\partial y^m}{\partial x^a}
 \frac{\partial x^b}{\partial y^m}
 T^a{}_b \notag\\
 &=\delta^b{}_aT^a{}_b
 =T^a{}_a.
\label{eq:trace-invariant}
\end{align}
У свёрнутой величины нет свободных индексов, и потому это скаляр.

\begin{figure}[H]
\centering
\includegraphics[width=\linewidth]{lecture_02_figure_20.png}
\caption{Преобразование смешанного тензора до свёртки верхнего и нижнего индексов.}
\lectureframe{01:07:45}
\end{figure}

\begin{figure}[H]
\centering
\includegraphics[width=\linewidth]{lecture_02_figure_21.png}
\caption{При свёртке якобианы сокращаются, оставляя один и тот же скалярный след в любой из двух систем координат.}
\lectureframe{01:10:30}
\end{figure}

\begin{classroomqa}
\audiencequestion{Остаются ли у результата индексы после свёртки?}
\lecturerresponse{Исчезает только свёртываемая пара. Свёртка единственного верхнего и нижнего индексов у \(T^m{}_n\) оставляет скаляр. В тензоре более высокого ранга все несвёрнутые индексы остаются свободными.
\lecturetimestamp{01:09:00}}
\end{classroomqa}

Например,
\begin{equation}
 A^{mn}{}_{pq}
 \longmapsto
 A^{mn}{}_{pn}
 =B^m{}_p
\label{eq:higher-rank-contraction}
\end{equation}
выполняет свёртку по \(n\) и уменьшает ранг на 2. Конкретные буквы — лишь немые обозначения; структурное правило состоит в суммировании по одному верхнему и одному нижнему индексу.

\begin{figure}[H]
\centering
\includegraphics[width=\linewidth]{lecture_02_figure_22.png}
\caption{Свёртка тензора более высокого ранга устраняет один верхний и один нижний индекс, сохраняя остальные свободные индексы.}
\lectureframe{01:12:35}
\end{figure}

\begin{classroomqa}
\audiencequestion{Почему нельзя таким же образом сворачивать два верхних индекса или два нижних?}
\lecturerresponse{Два верхних индекса дают два якобиана одной и той же направленности,
\(\partial y^m/\partial x^a\) и
\(\partial y^m/\partial x^b\). Они не перемножаются в тождество, поэтому предложенная сумма не подчиняется тензорному закону преобразования. Сначала можно с помощью метрики опустить один из индексов, после чего верхне--нижняя свёртка допустима; без этого дополнительного тензора операция не является инвариантной относительно координат. \lecturetimestamp{01:13:03}}
\end{classroomqa}

\begin{figure}[H]
\centering
\includegraphics[width=\linewidth]{lecture_02_figure_23.png}
\caption{Попытка свёртки индексов одной и той же валентности оставляет два якобиана одного направления и потому не сводится к тождеству.}
\lectureframe{01:14:30}
\end{figure}

\begin{classroomqa}
\audiencequestion{Является ли свёртка тензорной версией скалярного произведения?}
\lecturerresponse{Свёртка спаривает одно векторное место с одним ковекторным и является основным инвариантным спариванием. При наличии метрики два вектора можно превратить в такую пару, так что привычное скалярное произведение — это метрика, за которой следует свёртка. \lecturetimestamp{01:14:43}}
\end{classroomqa}

\begin{classroomqa}
\audiencequestion{Должны ли два сворачиваемых индекса пробегать один и тот же диапазон?}
\lecturerresponse{Да. Они относятся к двойственным базисам одного и того же \(d\)-мерного касательного пространства и кокасательного пространства. Иначе ни суммирование, ни тождество для обратных якобианов не были бы определены.
\lecturetimestamp{01:15:13}}
\end{classroomqa}

\begin{classroomqa}
\audiencequestion{Почему тензор, обращающийся в нуль в одной системе координат, обращается в нуль во всякой?}
\lecturerresponse{Каждая преобразованная компонента является линейной комбинацией исходных компонент. Все эти комбинации обращаются в нуль, когда каждая входная компонента равна нулю. Именно эта простая линейность делает тензорные уравнения независимыми от координат. \lecturetimestamp{01:15:31}}
\end{classroomqa}

\section{Метрика, задаваемая квадратом расстояния}

Теперь вернёмся к геометрии. Для бесконечно малого смещения с компонентами \(dx^m\) определим квадрат его длины как
\begin{equation}
 ds^2=g_{mn}(x)\,dx^m dx^n.
\label{eq:line-element}
\end{equation}
Это обобщённая теорема Пифагора. Метрика может меняться от точки к точке, а внедиагональные члены отражают неортогональность координатных направлений.

\begin{figure}[H]
\centering
\includegraphics[width=\linewidth]{lecture_02_figure_24.png}
\caption{Перечень операций над тензорами уступает место квадратичному определению бесконечно малого линейного элемента.}
\lectureframe{01:18:45}
\end{figure}

Вклад вносит только симметричная часть, поскольку
\(dx^m dx^n=dx^n dx^m\). Поэтому можно положить
\begin{equation}
 g_{mn}=g_{nm}.
\label{eq:metric-symmetry}
\end{equation}
В \(d\) измерениях симметричная метрика имеет
\begin{equation}
 \frac{d(d+1)}{2}
\label{eq:metric-component-count}
\end{equation}
независимых компонент: 10 в четырёх измерениях, 6 в трёх и 3 в двух.

\begin{figure}[H]
\centering
\includegraphics[width=\linewidth]{lecture_02_figure_25.png}
\caption{Линейный элемент и симметричный массив компонент, используемый для подсчёта независимых компонент метрики.}
\lectureframe{01:21:35}
\end{figure}

\begin{classroomqa}
\audiencequestion{Почему в четырёх измерениях у метрики 10 независимых компонент, а не 16?}
\lecturerresponse{Есть 4 диагональных элемента и 6 неупорядоченных внедиагональных пар. Эквивалентно,
\(4(4+1)/2=10\). Антисимметричная часть умножалась бы на симметричное произведение \(dx^m dx^n\) и потому взаимно уничтожалась бы. \lecturetimestamp{01:19:50}}
\end{classroomqa}

\begin{classroomqa}
\audiencequestion{Сколько независимых компонент в трёх измерениях?}
\lecturerresponse{Шесть: 3 диагональных элемента и 3 внедиагональные пары, или \(3(3+1)/2\). \lecturetimestamp{01:20:47}}
\end{classroomqa}

\begin{classroomqa}
\audiencequestion{А в двух измерениях?}
\lecturerresponse{Три: \(g_{11}\), \(g_{22}\) и единственная симметричная пара \(g_{12}=g_{21}\). \lecturetimestamp{01:22:28}}
\end{classroomqa}

\subsection{Почему метрика является тензором}

Квадрат расстояния — скаляр. Поэтому, записывая одно и то же смещение в двух системах координат, получаем
\begin{equation}
 g_{mn}(x)\,dx^m dx^n
 =
 g'_{pq}(y)\,dy^pdy^q.
\label{eq:length-invariance}
\end{equation}
Дважды используем
\begin{equation}
 dx^m=\frac{\partial x^m}{\partial y^p}\,dy^p
\end{equation}
:
\begin{align}
 ds^2
 &=
 g_{mn}(x)
 \frac{\partial x^m}{\partial y^p}
 \frac{\partial x^n}{\partial y^q}
 dy^pdy^q.
\label{eq:metric-transform-derivation}
\end{align}
Сравнение с выражением со штрихами даёт
\begin{equation}
 g'_{pq}(y)
 =
 \frac{\partial x^m}{\partial y^p}
 \frac{\partial x^n}{\partial y^q}
 g_{mn}(x).
\label{eq:metric-transformation}
\end{equation}
Это в точности закон преобразования ковариантного тензора второго ранга.

\begin{figure}[H]
\centering
\includegraphics[width=\linewidth]{lecture_02_figure_26.png}
\caption{Вывод закона преобразования метрики из инвариантности скалярного линейного элемента.}
\lectureframe{01:26:25}
\end{figure}

Этот вывод также объясняет два нижних индекса. Каждый нижний индекс метрики спаривается с одним верхним индексом, который несёт компонента смещения. При замене координат два обратных якобиана в
\eqref{eq:metric-transformation} сокращают два прямых якобиана дифференциалов, оставляя \(ds^2\) неизменным.

\section{Матрица метрики и её обратная}

В данной точке \(g_{mn}\) представляет собой симметричную матрицу:
\begin{equation}
 [g_{mn}]
 =
 \begin{pmatrix}
 g_{11} & g_{12} & \cdots \\
 g_{12} & g_{22} & \cdots \\
 \vdots & \vdots & \ddots
 \end{pmatrix}.
\label{eq:metric-matrix}
\end{equation}

\begin{figure}[H]
\centering
\includegraphics[width=\linewidth]{lecture_02_figure_27.png}
\caption{Метрика, записанная в виде симметричной матрицы после установления её тензорного закона преобразования.}
\lectureframe{01:28:50}
\end{figure}

Метрика должна быть невырожденной: у неё нет нулевого собственного значения, и потому она имеет обратную. Обозначим эту обратную через \(g^{mn}\).
В индексной записи,
\begin{equation}
 g_{mn}g^{np}=\delta_m{}^p.
\label{eq:inverse-metric}
\end{equation}
По повторяющемуся \(n\) ведётся суммирование, а \(m\) и \(p\) обозначают строку и столбец произведения.

\begin{figure}[H]
\centering
\includegraphics[width=\linewidth]{lecture_02_figure_28.png}
\caption{Обратная метрика, введённая посредством
\(g_{mn}g^{np}=\delta_m{}^p\).}
\lectureframe{01:31:40}
\end{figure}

\begin{classroomqa}
\audiencequestion{Где, собственно, живут векторы и тензоры в точке?}
\lecturerresponse{Векторы в точке образуют её касательное пространство
\(T_P M\); ковекторы образуют двойственное пространство \(T_P^*M\). Тензоры более высокого ранга принадлежат тензорным произведениям этих пространств. Тензорное поле гладко сопоставляет такой объект каждой точке. Метрика является билинейной формой на касательном пространстве, а её обратная — соответствующей билинейной формой на кокасательном пространстве. \lecturetimestamp{01:34:15}}
\end{classroomqa}

\begin{figure}[H]
\centering
\includegraphics[width=\linewidth]{lecture_02_figure_29.png}
\caption{Вектор, изображённый в касательном пространстве, прикреплённом к точке искривлённого многообразия.}
\lectureframe{01:35:45}
\end{figure}

\begin{classroomqa}
\audiencequestion{Относится ли
\(g_{mn}g^{np}=\delta_m{}^p\) только к случаю \(m=p\)?}
\lecturerresponse{Нет. Это полное матричное уравнение. Для каждой упорядоченной пары \((m,p)\) суммированное произведение равно 1, когда \(m=p\), и 0, когда \(m\ne p\). Все диагональные и внедиагональные уравнения присутствуют одновременно. \lecturetimestamp{01:36:50}}
\end{classroomqa}

\begin{figure}[H]
\centering
\includegraphics[width=\linewidth]{lecture_02_figure_30.png}
\caption{Обратное соотношение, интерпретированное как полная единичная матрица со смешанными индексами.}
\lectureframe{01:38:50}
\end{figure}

В двух измерениях это означает
\begin{equation}
 \begin{pmatrix}
 g_{11} & g_{12}\\
 g_{12} & g_{22}
 \end{pmatrix}
 \begin{pmatrix}
 g^{11} & g^{12}\\
 g^{12} & g^{22}
 \end{pmatrix}
 =
 \begin{pmatrix}
 1&0\\
 0&1
 \end{pmatrix}.
\label{eq:two-dimensional-inverse}
\end{equation}
Например,
\begin{align}
 g_{11}g^{11}+g_{12}g^{12}&=1,\\
 g_{11}g^{12}+g_{12}g^{22}&=0.
\end{align}
Другая строка даёт оставшиеся два уравнения.

\begin{figure}[H]
\centering
\includegraphics[width=\linewidth]{lecture_02_figure_31.png}
\caption{Тождество для обратной метрики, развёрнутое в явное двумерное матричное произведение.}
\lectureframe{01:40:00}
\end{figure}

\begin{classroomqa}
\audiencequestion{Обязательно ли метрика должна быть положительно определённой?}
\lecturerresponse{Для обычной римановой геометрии положительная определённость делает положительным квадрат длины любого ненулевого вектора. Вместо этого в пространстве-времени используется лоренцева метрика, у которой один знак противоположен трём остальным. Однако в обоих случаях здесь требуется невырожденность, а не положительность.
\lecturetimestamp{01:42:18}}
\end{classroomqa}

\begin{classroomqa}
\audiencequestion{Почему исключается нулевое собственное значение?}
\lecturerresponse{Нулевое собственное значение делает матрицу метрики вырожденной, так что \(g^{mn}\) не существует и метрика не может установить взаимно-однозначное соответствие между векторами и ковекторами. Метрика в обычном дифференциально-геометрическом смысле должна быть невырожденной. \lecturetimestamp{01:43:56}}
\end{classroomqa}

\begin{classroomqa}
\audiencequestion{Имеют ли 10 независимых компонент четырёхмерной метрики какое-либо отношение к 10-мерной теории струн?}
\lecturerresponse{Нет. Первое 10 — это просто
\(4(4+1)/2\), число элементов симметричной матрицы 4×4. Появление 10 измерений пространства-времени в теории суперструн имеет другое происхождение. \lecturetimestamp{01:44:32}}
\end{classroomqa}

\begin{classroomqa}
\audiencequestion{Если в лоренцевом пространстве-времени существуют ненулевые векторы нулевой длины, делает ли это метрику необратимой?}
\lecturerresponse{Нет. Изотропный вектор удовлетворяет
\(g_{mn}V^mV^n=0\), это взаимная компенсация внутри неопределённой квадратичной формы. Он не является вектором с нулевым собственным значением, удовлетворяющим
\(g_{mn}V^n=0\). Например, метрика Минковского имеет собственные значения
\(\pm1\) и вполне обратима, хотя и допускает светоподобные векторы. \lecturetimestamp{01:45:02}}
\end{classroomqa}

\section{Что дальше}

Теперь можно чётко сформулировать геометрическую программу. Метрика задаёт бесконечно малое расстояние, преобразуется ковариантно и имеет обратную. Её компоненты могут выглядеть сложными лишь потому, что сложны координаты. Чтобы решить, искривлена ли сама геометрия, мы должны сравнивать метрику в соседних точках независимым от координат образом.

Обычных производных недостаточно, потому что при дифференцировании тензора дифференцируются и зависящие от положения якобианы в его законе преобразования. Поэтому следующий шаг — ковариантное дифференцирование. Из производных метрики строят связность; из производных и произведений этой связности строят кривизну. Лишь тогда вопрос о согнутой странице получит свой окончательный внутренний ответ.
